%------------------------------------------------------------------------------%
\begin{question}
  Considerando a função
  \(
    f(x, y) = 4xy -x^4 -y^4
  \)

  a) Calcule o gradiente de $f$

  b) Calcule a hessiana de $f$

  c) Encontre todos os pontos críticos de $f$

  d) Classifique cada ponto crítico de $f$
\end{question}

%------------------------------------------------------------------------------%
\begin{resolution}{a -- Gradiente de $f$}
  \[
    \D{f}{x}
    \pause
    \;=\; \D{}{x}\left[4xy -x^4 -y^4\right]
    \pause
    \;=\; 4y - 4x^3
  \]

  \pause
  \[
    \D{f}{y}
    \pause
    \;=\;  \D{}{y}\left[4xy -x^4 -y^4\right]
    \pause
    \;=\;  4x -4y^3
  \]

  \pause
  \[
    \text{então}\qquad
    \nabla f(x, y) \;=\; \vetor{4y - 4x^3}{4x -4y^3}
  \]
\end{resolution}

%------------------------------------------------------------------------------%
\begin{resolution}{b -- Hessiana de $f$}
  \[
    \DS{f}{x}
    \pause
    \;=\; \D{}{x}\left[4y - 4x^3\right]
    \pause
    \;=\; -12x^2
  \]

  \pause
  \[
    \DS{f}{y}
    \pause
    \;=\; \D{}{y}\left[4x -4y^3\right]
    \pause
    \;=\; -12y^2
  \]

  \pause
  \[
    \DM{f}{x}{y}
    \pause
    \;=\; \D{}{x}\D{f}{y}
    \pause
    \;=\; \D{}{x}\left[4x -4y^3\right]
    \pause
    \;=\; 4
  \]
  \pause
  \[
    \text{então}\qquad
    H(x, y)
    \;=\; \left(\begin{array}{cc}
      -12x^2 & 4      \\
      4      & -12y^2
    \end{array}\right)
  \]
\end{resolution}

%------------------------------------------------------------------------------%
\begin{resolution}{c -- Pontos críticos de $f$}
  A função é um polinômio, então possui derivadas em todos os pontos do plano

  \pause
  Assim os pontos críticos são apenas os pontos onde \(\nabla f = 0\)

  \pause
  \[
    4y - 4x^3 = 0
    \qquad\text{e}\qquad
    4x -4y^3 = 0
  \]

  \pause
  ou
  \[
    y - x^3 = 0
    \qquad\text{e}\qquad
    x - y^3 = 0
  \]
\end{resolution}

%------------------------------------------------------------------------------%
\begin{resolution}{c -- Resolvendo o sistema}
\begin{columns}[t]
\column{0.45\linewidth}
  \onslide<+->{Isolando $y$ na primeira equação}
  \begin{align*}
    \onslide<+->{y - x^3 & = 0   \\}
    \onslide<+->{y       & = x^3}
  \end{align*}
\column{0.55\linewidth}
  \onslide<+->{e substituindo na segunda}
  \begin{align*}
    \onslide<+->{x - y^3                & = 0 \\}
    \onslide<+->{x - \left(x^3\right)^3 & = 0 \\}
    \onslide<+->{x -x^9                 & = 0 \\}
    \onslide<+->{x\left(1 - x^8\right)  & = 0}
  \end{align*}
  \onslide<+->{As soluções dessa equação são\\$x=0$, $x=1$ ou $x-1$}
\end{columns}
\end{resolution}

%------------------------------------------------------------------------------%
\begin{resolution}{c -- Resolvendo o sistema}
  Como \(y=x^3\)

  \pause
  Se \(x=0\) temos \(y=0\)

  \pause
  Se \(x=1\) temos \(y=1\)

  \pause
  Se \(x=-1\) temos \(y=-1\)

  \pause
  \bigskip
  Os pontos críticos são
  \[
    (x_1, y_1) = (0, 0),
    \qquad
    (x_2, y_2) = (1, 1)
    \quad\text{e}\quad
    (x_3, y_3) = (-1, -1)
  \]
\end{resolution}

%------------------------------------------------------------------------------%
\begin{resolution}{d -- Classificando o ponto \((x_1, y_1) = (0, 0)\)}
\begin{columns}[t]
\column{0.45\linewidth}
  \[
    f_{xx}(0, 0)
    \pause
    \;=\; \left(-12x^2\right) \evalat{(0, 0)}{}
    \pause
    =\; 0
  \]
  \[
    \pause
    f_{yy}(0, 0)
    \pause
    \;=\; \left(-12y^2\right) \evalat{(0, 0)}{}
    \pause
    =\; 0
  \]
  \[
    \pause
    f_{xy}(0, 0)
    \pause
    \;=\; \left(4\right) \evalat{(0, 0)}{}
    \pause
    =\; 4
  \]
\column{0.55\linewidth}
  \pause
  \begin{align*}
    \onslide<+->{D_1}
    \onslide<+->{& = \det H(0, 0) \\}
    \onslide<+->{& = f_{xx}(0, 0)f_{yy}(0, 0) - f^2_{xy}(0, 0) \\}
    \onslide<+->{& = 0 \times 0 - 4^2 \\}
    \onslide<+->{& = -16 }
    \onslide<+->{< 0}
  \end{align*}
\end{columns}

  \bigskip
  \onslide<+->{O ponto $(0, 0)$ é um \emph{ponto de sela}}
\end{resolution}

%------------------------------------------------------------------------------%
\begin{resolution}{d -- Classificando o ponto \((x_2, y_2) = (1, 1)\)}
\begin{columns}[t]
\column{0.45\linewidth}
  \[
    f_{xx}(1, 1)
    \pause
    \;=\; \left(-12x^2\right) \evalat{(1, 1)}{}
    \pause
    =\; -12
  \]
  \[
    \pause
    f_{yy}(1, 1)
    \pause
    \;=\; \left(-12y^2\right) \evalat{(1, 1)}{}
    \pause
    =\; -12
  \]
  \[
    \pause
    f_{xy}(1, 1)
    \pause
    \;=\; \left(-12x^2\right) \evalat{(1, 1)}{}
    \pause
    =\; 4
  \]
\column{0.55\linewidth}
  \pause
  \begin{align*}
    \onslide<+->{D_2}
    \onslide<+->{& = \det H(1, 1) \\}
    \onslide<+->{& = f_{xx}(1, 1)f_{yy}(1, 1) - f^2_{xy}(1, 1) \\}
    \onslide<+->{& = (-12)(-12) - 4^2 \\}
    \onslide<+->{& = 144 -16 \\}
    \onslide<+->{& = 128 }
    \onslide<+->{> 0}
  \end{align*}
\end{columns}

  \bigskip
  \onslide<+->{O ponto $(1, 1)$ é um extremo local}

  \onslide<+->{Como \(f_{xx}(1, 1) = -12 < 0\)}
  \onslide<+->{o ponto é um \emph{máximo local}}
\end{resolution}

%------------------------------------------------------------------------------%
\begin{resolution}{d -- Classificando o ponto \((x_3, y_3) = (-1, -1)\)}
\begin{columns}[t]
\column{0.5\linewidth}
  \[
    f_{xx}(-1, -1)
    \pause
    = \left(-12x^2\right) \evalat{(-1, -1)}{}
    \pause
    \!\!= -12
  \]
  \[
    \pause
    f_{yy}(-1, -1)
    \pause
    = \left(-12x^2\right) \evalat{(-1, -1)}{}
    \pause
    \!\!= -12
  \]
  \[
    \pause
    f_{xy}(-1, -1)
    \pause
    = \left(-12x^2\right) \evalat{(-1, -1)}{}
    \pause
    \!\!= 4
  \]
\column{0.5\linewidth}
  \pause
  \begin{align*}
    \!
    \onslide<+->{D_3}
    \onslide<+->{& = \det H(-1, -1) \\}
    \onslide<+->{& = f_{xx}(-1, -1)f_{yy}(-1, -1) - f^2_{xy}(-1, -1) \\}
    \onslide<+->{& = (-12)(-12) - 4^2 \\}
    \onslide<+->{& = 144 -16 \\}
    \onslide<+->{& = 128 }
    \onslide<+->{> 0}
  \end{align*}
\end{columns}

  \bigskip
  \onslide<+->{O ponto $(-1, -1)$ é um extremo local}

  \onslide<+->{Como \(f_{xx}(-1, -1) = -12 < 0\)}
  \onslide<+->{o ponto é um \emph{máximo local}}
\end{resolution}

%------------------------------------------------------------------------------%
